\scanpage{13}
\mspage{10}

Cette relation d'ordre a un plus petit élément, $x_0$, et un plus grand élément, $x_1$. Elle permet de reconstituer la structure de modèle de $1$-forme. En effet, l'ens des segments \deleted{est} \inserted{s'identifie à} \unsure{celui} des couples $\{x,y\}$ tels que $x<y$, et pour un tel segment $I_{x,y}$, on a
\[
\tilde I_{xy}=\{z\in L\mid x\leq z\leq y\}
\]
\[
\delta I_{x,y}=\{x,y\}.
\]

Cette relation d'ordre est ``infiniment divisible'' i.e. si $x,y$ sont tels que $x<y$, $\exists z$ tel que $x<z<y$.

\inserted{Elle est \unsure{loc. filtrante croissante} et filtrante décroissante, i.e. si $x<y_1,y_2$ $\exists y$ tel que $x<y\leq y_1,y_2$, et dualement$^\ast$.}

\margin{$^\ast$ Mais pour ceci il faut \unsure{compléter} F3) en posant : Si $I$ est un segment, $x\in\delta I$, $y,z\in I^\circ$, alors $I_{x,y}^\circ\cap I_{x,z}^\circ\ne\varnothing$ ; $\exists u\in I$ tel que $I_{x,\unsure{u|y}}\subset I_{x,y}\cap I_{x,z}$.}

Inversement, une relation d'ordre sur $L$, divisible \inserted{(et loc. bifiltr.)}, ayant un plus petit élément \deleted{\unsure{et}} $x_0$ et un plus grand élément $x_1$, définit (avec la \unsure{notation|notion} précédente de segment \inserted{$I_{x,y}$} et \deleted{\unsure{la}} la description de $\tilde I_{xy}$ et $\delta I_{xy}$)

\mspage{12}

\rotatebox{90}{\unsure{EXTRA STRONG}}
\rotatebox{90}{\unsure{EXTRA STRONG}}
